\documentclass[11pt]{article}

% --------------------
% Packages
% --------------------
\usepackage[a4paper,margin=1in]{geometry}
\usepackage{setspace}
\usepackage{hyperref}
\usepackage{enumitem}
\usepackage{titlesec}
\usepackage{hyperref}
\usepackage{tcolorbox}
\tcbuselibrary{documentation, skins}
\usepackage{longtable}
\setstretch{1.1}

% --------------------
% Formatting tweaks
% --------------------
\titleformat{\section}
  {\normalfont\Large\bfseries}{\thesection}{1em}{}

\titleformat{\subsection}
  {\normalfont\large\bfseries}{\thesubsection}{1em}{}

\titleformat{\subsubsection}
  {\normalfont\normalsize\bfseries}{\thesubsubsection}{1em}{}

\setlist[itemize]{noitemsep, topsep=4pt}

% --------------------
% Metadata
% --------------------
\title{Lab Manual}
\author{Roars Lab}
\date{Last updated: \today}

\begin{document}
\maketitle
\tableofcontents
\newpage

% ============================================================
\section{About this Manual}

\subsection{Purpose of this manual}
This manual serves as a guide for prospective and current members of \href{https://roars.dev}{Roars Lab}. It outlines the lab's policies, procedures, and expectations to ensure a productive and collaborative research environment.

\subsection{Updating the manual}
This document is a living document in Roars Lab's Github repo at

\begin{tcolorbox}[enhanced]
    \centering
\url{https://code.roars.dev/phd-cs-us/lab-manual.tex}.    
\end{tcolorbox}

Lab members, who all have write access, are encouraged to provide feedback and suggest changes to improve its clarity and usefulness.



% ============================================================
\section{About the Lab}

\subsection{Lab mission and focus}
Research goals, scope, and ambitions.

\subsection{Lab philosophy and values}
Open science, rigor, transparency, reproducibility, etc.

\subsection{Lab history}
Founding, major changes, milestones.

\subsection{Lab space}
Physical description of the lab.

\subsection{Contact information}
Address, email, website.


\begin{itemize}
    \item Lab website: \url{https://roars.dev}
    \item Email: \href{mailto:info@roars.dev}{info@roars.dev}
    \item Address
    \begin{tcolorbox}
George Mason University\\
Nguyen Engineering Building  \#4430\\
4400 University Drive\\
Fairfax, VA 22030
    \end{tcolorbox}
\end{itemize}


\subsection{Facilities}
The lab provides:
\begin{itemize}
  \item High-end desktops, laptops, and peripherals
  \item Shared GPU and compute servers
  \item Cloud and storage resources as needed
\end{itemize}

\subsection{Ethical code of conduct}

\subsubsection{Ethical principles}
Lab members are expected to treat each other with respect, act professionally,
and take responsibility for their work.

\subsubsection{Diversity statement}
Roars Lab values diversity in background, identity, and perspective.
We aim to provide an inclusive and respectful environment where people are judged
by the quality of their ideas and work.

\subsubsection{Institutional rules}
Lab members must follow GMU and department-level codes of conduct and policies.

% ============================================================
\section{Becoming a Lab Member}

\subsection{Online accounts}
Lab members are expected to maintain accounts on:
\begin{itemize}
  \item GitHub
  \item Discord (primary lab communication)
  \item Overleaf
  \item Google Scholar
  \item ORCID (recommended)
\end{itemize}

\subsection{Skills updating}
Members are expected to continuously update their skills.
Typical skills include programming, formal reasoning, writing, and tool building.
Training is largely self-directed.

\subsection{Required readings}
New members should read core papers relevant to their area and prior work from the lab.
Specific reading lists are project-dependent.

\subsection{Internal resources}
\begin{itemize}
  \item Lab servers
  \item Shared repositories
  \item Internal wikis and documentation
\end{itemize}

\subsection{External resources}
Institutional resources such as libraries, IT support, and administrative staff are available.

% ============================================================
\section{Being a Lab Member}

\subsection{Roles and Duties}

\subsubsection{Roles}
Roles exist mainly for clarity, not hierarchy.

\subsubsection{Full members}
PI, PhD students, MS students, and undergraduate researchers.

\subsubsection{Associated members}
Visitors and collaborators working on shared projects.

\subsubsection{Non-scientific members}
Administrative staff and coordinators as applicable.

\subsubsection{Duties and responsibilities}
All members are expected to:
\begin{itemize}
  \item Own their projects
  \item Help others when possible
  \item Contribute to a healthy lab environment
\end{itemize}

\subsection{Career Development}

\subsubsection{Mentoring}
Mentoring is individualized.
Expectations increase with seniority.
Letters of recommendation require advance notice and evidence of work.

\subsection{How We Do Things}

\subsubsection{Time management}

\paragraph{Annual structure}
The lab follows the academic calendar, with flexibility outside major deadlines.

\paragraph{Leave and holidays}
Take time off when needed.
Outside deadlines, no approval is required for short breaks.

\paragraph{Lab hours}
There are no fixed working hours.
Productivity matters more than presence.

\paragraph{Absences}
Extended absences should be communicated.

\paragraph{Meetings}
Weekly lab meetings are standard.
Everyone speaks; no one hides.

\paragraph{Deadlines}
Deadlines are taken seriously.
Work intensifies near submissions and relaxes afterward.

\paragraph{Remote work}
Remote work is acceptable as long as progress is made.

\paragraph{Calendar usage}
Use calendars sensibly for meetings and deadlines.

\paragraph{Social events}
Social events happen organically.
No attendance is mandatory.

\subsection{Communication}

\subsubsection{Conflict resolution}
Conflicts should be addressed early and directly.
The PI will step in when needed.

\subsubsection{Internal communication}
Discord is primary.
Email is for official matters.

\subsubsection{External communication}
Public communication should reflect well on the lab.

\subsubsection{Collaborators}
External collaborations should be transparent.
When unsure, keep the PI in the loop.

\subsubsection{Visiting researchers}
Visits require PI approval and planning.

\subsubsection{Lab vocabulary}
Shared technical language emerges naturally.
Ask when something is unclear.

\subsubsection{Life in the lab}
Be considerate about noise, space, and shared resources.

\subsubsection{Guests}
Guests require permission and must follow safety rules.

% ============================================================
\section{Dealing with Mistakes}

\subsection{General principles}
Mistakes happen.
Report them early.
Fix what can be fixed and learn from them.

\subsection{Confidentiality breaches}
Report immediately to the PI.

\subsection{Troubleshooting}
Document failures.
Reproduce issues.
Do not hide problems.

% ============================================================
\section{Maintaining Health and Safety}

\subsection{Health and safety}
Follow institutional safety procedures.
Seek help when sick or overwhelmed.

\subsection{Healthy lab life}
Burnout is not a virtue.
Support each other.

% ============================================================
\section{Leaving the Lab}

\subsection{Contacts and timeline}
Discuss plans early.

\subsection{Access revocation}
Accounts and access will be revoked after departure.

\subsection{Letters of recommendation}
Request well in advance and provide materials.

\subsection{Future steps}
The PI will assist with career transitions when possible.

\subsection{Keeping in touch}
Former members are encouraged to stay in contact.

% ============================================================
\section{Preparing a Study}

\subsection{Initiating a Study}

\subsubsection{Starting a new project}
Projects require PI approval.

\subsubsection{Project team}
Teams are small and flexible.

\subsubsection{Project logging}
All projects must have a documented home.

\subsection{Feasibility and Validity}

\subsubsection{Consultation}
Do literature reviews before committing.

\subsubsection{Cross-checking}
Test assumptions early.

\subsubsection{Pilot testing}
Pilot when appropriate.

\subsection{Operational Steps}

\subsubsection{Preregistration}
Used when appropriate.

\subsubsection{Approvals}

\paragraph{IRB approval}
Required for human-subject research.

\paragraph{Other approvals}
As needed.

\subsubsection{Data management plan}
Required for funded projects.

\subsubsection{Time management plan}
Milestones should be explicit.

\subsubsection{Materials and equipment}
Plan ahead.

\subsubsection{Experiment implementation}
Use lab-standard tools when possible.

\subsubsection{Testing space}
Coordinate shared resources.

% ============================================================
\section{Conducting a Study}

\subsection{Recruitment}
Follow approved protocols.

\subsection{Equipment use}
Training may be required.

\subsection{Data collection}
Follow documented procedures.

\subsection{Protocol compliance}
Checklists are encouraged.

\subsection{Documentation}
Log decisions and deviations.

\subsection{Closing data collection}
Archive and summarize.

% ============================================================
\section{Data and File Management}

\subsection{Data processing and analysis}
Analysis must be reproducible.

\subsection{Quality control}
Code should be readable and maintainable.

\subsection{Metadata}
Use codebooks and documentation.

\subsection{Standards}

\subsubsection{Dataset structure}
Prefer tidy, well-structured formats.

\subsubsection{Folder structure}
Separate raw and processed data.

\subsubsection{Naming conventions}
Be consistent.

\subsubsection{File formats}
Use open formats when possible.

\subsubsection{Version control}
Git is strongly encouraged.

\subsubsection{Archiving}
Archive all final artifacts.

\subsubsection{Data privacy}
Protect sensitive data.

\subsubsection{Sharing and access}
Follow open science norms when possible.

\subsubsection{Intellectual property}
Follow institutional IP policies.

% ============================================================
\section{Communicating Results}

\subsection{Reporting standards}
Follow community guidelines.

\subsection{Plotting and visualization}
Clarity over decoration.

\subsection{Writing papers}
Overleaf is standard.

\subsection{Proofreading}
Peer review within the lab is encouraged.

\subsection{Choosing venues}
Quality over quantity.

\subsection{Accreditation}

\subsubsection{Contributorship}
Decide early and document.

\subsubsection{Acknowledgments}
Acknowledge funding and help.

\subsection{Open access}
Use preprints when appropriate.

\subsection{Posters}
Use clean, readable templates.

\subsection{Talks}
Practice matters.

\subsection{Public engagement}
Optional, but encouraged.

% ============================================================
\section{Funding and Budgeting}

\subsection{Lab funds}
The lab covers essential research expenses.

\subsection{Student funding}
Students are supported via grants, fellowships, or TA-ships.

\subsection{Grants and administration}
The PI handles grant administration with student input when needed.

% ============================================================
\section*{Contributors}
List of contributors to the lab manual.

\end{document}
